\documentclass[11pt,a4paper,fontset=windows]{ctexart}
% ===== 公共导言区：高等代数（一）期中模拟卷 =====
\usepackage{amsmath,amssymb}
\usepackage{unicode-math}
\setmathfont{TeX Gyre Termes Math}
\usepackage[margin=2.05cm,top=2.35cm,bottom=2.0cm,headsep=11pt,includefoot]{geometry}
\usepackage{fancyhdr}
\usepackage{lastpage}
\usepackage{enumitem}
\usepackage{array}
\usepackage{bm}
\usepackage{xcolor}
\usepackage{zhnumber}
\usepackage{needspace}

\definecolor{hdgray}{gray}{0.32}
\definecolor{lngray}{gray}{0.5}
\definecolor{notebg}{gray}{0.955}

\pagestyle{fancy}
\fancyhf{}
\renewcommand{\headrulewidth}{0.4pt}
\renewcommand{\footrulewidth}{0pt}
\fancyhead[L]{\small\color{hdgray}\songti 高等代数（一）\ 期中模拟试卷（\thepp）}
\fancyhead[R]{\small\color{hdgray}\songti 考试时间：120 分钟\quad 满分：100 分}
\fancyfoot[C]{\small\color{lngray}\songti 第 \thepage\ 页\quad 共 \pageref{LastPage}\ 页}

\setlength{\parindent}{0pt}
\setlength{\parskip}{2pt}
\linespread{1.16}

% 题目列表 1—8 连续编号
\newlist{qs}{enumerate}{1}
\setlist[qs]{label=\textbf{\arabic*.},ref=\arabic*,leftmargin=2.55em,
  labelsep=0.55em,itemsep=23pt,topsep=6pt,parsep=0pt}
\newlist{note}{enumerate}{1}
\setlist[note]{label=(\arabic*),leftmargin=1.75em,itemsep=1.5pt,topsep=2.5pt,parsep=0pt}
\newlist{sub}{enumerate}{1}
\setlist[sub]{label=(\arabic*),leftmargin=1.9em,itemsep=3pt,topsep=3pt,parsep=0pt}

\newcommand{\secA}[1]{\Needspace{9\baselineskip}\par\vspace{10pt}{\large\bfseries\songti #1}\par\vspace{3pt}\nopagebreak}
\newcommand{\blank}[1][2.6cm]{\underline{\hspace{#1}}}
\newcounter{pp}
\renewcommand{\thepp}{\chinese{pp}}

\setcounter{pp}{4}
\begin{document}

\begin{center}
{\LARGE\bfseries\songti 高等代数（一）期中模拟试卷（四）}\\[7pt]
{\normalsize\songti 适用范围：第一章\ 多项式（\S1\,--\,\S9）；第二章\ 行列式（\S2.1\,--\,\S2.7）}\\[3pt]
{\small\color{hdgray}\songti 多元多项式、对称多项式不在考查范围}
\end{center}
\vspace{-4pt}

\begin{center}
\begin{tabular}{@{}p{0.335\textwidth}p{0.335\textwidth}p{0.26\textwidth}@{}}
\songti 姓名：\blank[2.5cm] & \songti 学号：\blank[2.5cm] & \songti 得分：\blank[1.9cm]
\end{tabular}
\end{center}

\vspace{2pt}
\hrule height 0.5pt
\begin{center}
\setlength{\fboxsep}{6pt}%
\colorbox{notebg}{\parbox{0.94\textwidth}{\small\songti
\textbf{考生须知：}
\begin{note}
\item 本试卷共 8 题，满分 100 分，考试时间 120 分钟。
\item 计算题须写出必要的计算过程，证明题须给出完整的推理步骤；只写答案不给分。
\item 所有答案一律写在答题纸上，写在试卷上无效。
\end{note}}}
\end{center}
\hrule height 0.5pt
\vspace{4pt}

\secA{一、计算题（每小题 10 分，共 40 分）}
\begin{qs}
\item 计算 $n$ 阶行列式
\[
D_n=\begin{vmatrix}
0&1&2&3&\cdots&n-1\\
1&0&1&2&\cdots&n-2\\
2&1&0&1&\cdots&n-3\\
3&2&1&0&\cdots&n-4\\
\vdots&\vdots&\vdots&\vdots&\ddots&\vdots\\
n-1&n-2&n-3&n-4&\cdots&0
\end{vmatrix}.
\]

\item 设 $a_1a_2\cdots a_n\neq 0$。计算 $n+1$ 阶行列式
\[
D_{n+1}=\begin{vmatrix}
a_0&b_1&b_2&\cdots&b_n\\
c_1&a_1&0&\cdots&0\\
c_2&0&a_2&\cdots&0\\
\vdots&\vdots&&\ddots&\vdots\\
c_n&0&\cdots&0&a_n
\end{vmatrix}.
\]

\item 设 $\lambda\in\mathbb{C}$。求使多项式
\[
f(x)=x^{4}-4x+\lambda
\]
有重根的全部 $\lambda$，并对每个 $\lambda$ 写出相应的重根。

\item 设
\[
f(x)=2x^{3}-2x^{2}-x+1,\qquad g(x)=x^{2}-1 .
\]
求 $(f(x),\,g(x))$，并求 $u(x),v(x)\in\mathbb{Q}[x]$，使 $u(x)f(x)+v(x)g(x)=(f(x),g(x))$。
\end{qs}

\secA{二、证明题（每小题 15 分，共 60 分）}
\begin{qs}[resume*]
\item 设 $c$ 为非零常数，$f(x)\in P[x]$，且
\[
f(x-c)=f(x).
\]
证明：$f(x)$ 必为常数。

\item 设 $f(x),g(x)\in P[x]$，且 $(f(x),g(x))=1$。证明：
\begin{sub}
\item 对任意正整数 $m,n$，都有 $\bigl(f^{m}(x),\,g^{n}(x)\bigr)=1$；
\item $\bigl(f(x)g(x),\;f(x)+g(x)\bigr)=1$。
\end{sub}

\item 设
\[
f(x)=a_nx^{n}+a_{n-1}x^{n-1}+\cdots+a_0\in\mathbb{Z}[x],\qquad a_na_0\neq 0 .
\]
\begin{sub}
\item 证明：若 $\dfrac{p}{q}$（$p,q$ 为互素的整数，$q>0$）是 $f(x)$ 的有理根，
      则 $p\mid a_0$，$q\mid a_n$；
\item 由此证明：首一整系数多项式的有理根必为整数。
\end{sub}

\item 设 $A$ 为 $m$ 阶方阵，$C$ 为 $n$ 阶方阵，$B$ 为 $m\times n$ 矩阵。证明：
\[
\begin{vmatrix}
A&B\\[2pt]
0&C
\end{vmatrix}=|A|\cdot|C| .
\]
\end{qs}

\end{document}
